\chapter{Ortalamaların Karşılaştırılması: t Testleri}
\label{ch:power-one-sample}
\label{ch:two-sample}
\index{t testi@$t$ testi}

\begin{hedefler}
Bu bölümün sonunda okuyucu:
\begin{itemize}
  \item araştırma tasarımına göre tek örneklem, bağımsız örneklem ve eşleştirilmiş örneklem $t$ testini seçebilecek,
  \item her test için hedef parametreyi ve hipotezleri açıkça yazabilecek,
  \item test istatistiği, serbestlik derecesi, $p$-değeri ve güven aralığını birlikte yorumlayabilecek,
  \item bağımsız gruplarda klasik ve Welch $t$ testini ayırabilecek,
  \item eşleştirilmiş testte analizin neden fark puanları üzerinden yürütüldüğünü açıklayabilecek,
  \item bağımsızlık, dağılım biçimi, aykırı değer ve varyans yapısını uygun biçimde değerlendirebilecek,
  \item Cohen'in $d$ değeri ve ham ortalama farkıyla uygulamalı önemi tartışabilecek,
  \item Python ve SPSS çıktılarından eksiksiz bir bulgu raporu oluşturabilecektir.
\end{itemize}
\end{hedefler}

\begin{hazirlik}
Bir grubun ön-test ve son-test puanlarını karşılaştırmak ile iki farklı grubun
puanlarını karşılaştırmak neden aynı tasarım değildir? Gözlemlerin birbirine
bağlı veya bağımsız olması açısından düşünün.
\end{hazirlik}

\section{PDR vakası: Üç farklı ortalama sorusu}

Bir üniversitenin psikolojik danışma merkezi aynı yalnızlık ölçeğiyle üç
araştırma sorusu sormaktadır:

\begin{enumerate}
  \item Birinci sınıf öğrencilerinin ortalama puanı norm değeri 50'den farklı mıdır?
  \item Grup danışmanlığına katılan öğrenciler ile karşılaştırma grubundaki farklı öğrencilerin son-test ortalamaları farklı mıdır?
  \item Programa katılan aynı öğrencilerin ön-test ve son-test ortalamaları değişmiş midir?
\end{enumerate}

Üç soru da ortalamaya ilişkindir; fakat veri yapıları farklıdır. İlk soruda
bir örneklem ortalaması sabit bir değerle, ikinci soruda farklı kişilerin iki
ortalaması, üçüncü soruda ise aynı kişilerin iki ölçümü karşılaştırılır. Test
seçimi yalnız değişken sayısına değil, gözlemlerin nasıl üretildiğine dayanır.

\begin{figure}[htbp]
\centering
\resizebox{0.98\linewidth}{!}{%
\begin{tikzpicture}[alt={Veri yapısına göre tek örneklem, bağımsız örneklem ve eşleştirilmiş örneklem t testi seçim şeması},node distance=0cm]
  \node[draw=PrismBlue,fill=PrismLight,rounded corners,align=center,minimum width=3.3cm,minimum height=0.8cm] (q) at (0,2.4) {Karşılaştırmanın veri yapısı nedir?};
  \node[draw=PrismBlue!70,fill=white,rounded corners,align=center,minimum width=3.0cm,minimum height=0.8cm] (one) at (-4,0.5) {Bir grup\\sabit değer};
  \node[draw=PrismBlue!70,fill=white,rounded corners,align=center,minimum width=3.0cm,minimum height=0.8cm] (ind) at (0,0.5) {İki farklı\\bağımsız grup};
  \node[draw=PrismBlue!70,fill=white,rounded corners,align=center,minimum width=3.0cm,minimum height=0.8cm] (pair) at (4,0.5) {Aynı veya\\eşleştirilmiş birimler};
  \node[font=\small\bfseries,text=PrismBlue] at (-4,-0.7) {Tek örneklem t};
  \node[font=\small\bfseries,text=PrismBlue] at (0,-0.7) {Welch / bağımsız t};
  \node[font=\small\bfseries,text=PrismBlue] at (4,-0.7) {Eşleştirilmiş t};
  \draw[->,thick,PrismBlue] (q)--(one);
  \draw[->,thick,PrismBlue] (q)--(ind);
  \draw[->,thick,PrismBlue] (q)--(pair);
\end{tikzpicture}%
}
\caption{Araştırma tasarımına göre temel $t$ testi seçimi.}
\end{figure}

\begin{mistakebox}
Aynı öğrencilerin ön-test ve son-test ölçümlerini iki bağımsız grup gibi
analiz etmek eşleşme bilgisini kaybeder ve standart hatayı yanlış hesaplar.
Tersine, birbirinden farklı öğrencileri yalnız grup büyüklükleri eşit diye
eşleştirilmiş kabul etmek de geçersizdir.
\end{mistakebox}

\section{t testlerinin ortak mantığı}
\index{t testi@$t$ testi!test istatistiği}
\index{serbestlik derecesi}

Bütün $t$ testleri gözlenen bir tahmini sıfır hipotezindeki değerle
karşılaştırır \cite{casella2002,wasserman2004}:
\[
t=\frac{\text{gözlenen tahmin}-\text{$H_0$ altındaki değer}}
        {\text{tahminin standart hatası}}.
\]

Pay araştırmanın ilgilendiği farkı, payda ise bu farkın örnekleme
değişkenliğini gösterir. Mutlak $t$ değeri büyüdükçe gözlenen sonuç $H_0$
değerinden standart hata birimleriyle daha uzakta bulunur. İşaret farkın
yönünü gösterir; kanıtın derecesi çift yönlü testte mutlak değer üzerinden
değerlendirilir.

\begin{table}[htbp]
\centering
\caption{Üç $t$ testinin hedefleri ve analiz birimleri.}
\begin{tabular}{@{}p{0.22\textwidth}p{0.25\textwidth}p{0.31\textwidth}@{}}
\toprule
Test & Hedef parametre & Analize giren temel değer \\
\midrule
Tek örneklem & $\mu-\mu_0$ & Her katılımcının ham puanı \\
Bağımsız örneklem & $\mu_1-\mu_2$ & İki farklı gruptaki ham puanlar \\
Eşleştirilmiş örneklem & $\mu_D$ & Her eşleşme için fark puanı $D_i$ \\
\bottomrule
\end{tabular}
\end{table}

Testin sonunda yalnız $p$-değeri değil, hedef farkın nokta tahmini ve güven
aralığı da raporlanır. Böylece farkın yönü, olası büyüklüğü ve tahmin
hassasiyeti görünür olur.

\section{Tek örneklem t testi}
\index{t testi@$t$ testi!tek örneklem}

Tek örneklem $t$ testi, bir evren ortalamasını önceden tanımlanmış bir referans
değerle karşılaştırır. Hipotezler
\[
H_0:\mu=\mu_0,
\qquad
H_1:\mu\neq\mu_0
\]
ve test istatistiği
\[
t=\frac{\bar X-\mu_0}{s/\sqrt n},
\qquad sd=n-1
\]
biçimindedir.

Referans değer araştırma verisinden keyfî olarak seçilmemelidir. Ölçeğin norm
ortalaması, program hedefi veya kuramsal eşik olabilir; kaynağı ve güncelliği
açıklanmalıdır. Eski ya da farklı bir norm grubunun ortalaması, yeni öğrenci
evreni için anlamlı bir karşılaştırma oluşturmayabilir.

\subsection{Adım adım tek örneklem örneği}

Yirmi beş öğrencinin akademik kaygı puanı ortalaması 54,8 ve standart sapması
8 olsun. Güncel ve uygun norm ortalaması 50 ile çift yönlü karşılaştırılacaktır.

\adimbaslik{1. Hipotezler}
\[
H_0:\mu=50,
\qquad H_1:\mu\neq50.
\]

\adimbaslik{2. Standart hata ve test istatistiği}
\[
SE(\bar X)=\frac{8}{\sqrt{25}}=1{,}6,
\qquad
t=\frac{54{,}8-50}{1{,}6}=3{,}00.
\]

Serbestlik derecesi 24'tür ve çift yönlü $p$-değeri yaklaşık 0,006'dır.

\adimbaslik{3. Güven aralığı}

Yüzde 95 güven düzeyinde $t^*\approx2{,}064$ olduğundan
\[
54{,}8\pm2{,}064(1{,}6)
=[51{,}50;\ 58{,}10]
\]
elde edilir. Aralık 50'yi dışlar ve test kararıyla uyumludur.

\adimbaslik{4. Etki büyüklüğü}

Tek örneklem için standartlaştırılmış fark
\[
d=\frac{\bar x-\mu_0}{s}
=\frac{4{,}8}{8}=0{,}60
\]
olur. Sonuç norm değerinden istatistiksel olarak belirgin ve yaklaşık 0,60
standart sapma yüksek bir ortalamaya işaret eder. Bunun klinik anlamı ölçeğin
yorum kılavuzu ve PDR bağlamıyla ayrıca değerlendirilmelidir.

\begin{researchbox}
Bir grubun ortalamasının klinik eşikten yüksek olması, gruptaki bütün
öğrencilerin eşik üzerinde olduğu anlamına gelmez. Ortalama karşılaştırması
bireysel sınıflandırma yapmaz; puan dağılımı ve eşik üzerindeki öğrenci oranı
ayrıca incelenmelidir.
\end{researchbox}

\section{Bağımsız örneklem t testi}
\index{t testi@$t$ testi!bağımsız örneklem}
\index{bağımsız gruplar}

Bağımsız örneklem testinde iki gruptaki kişiler farklıdır ve bir kişinin
puanı diğer gruptaki kişinin puanını belirlemez. Hedef parametre
$\mu_1-\mu_2$ farkıdır:
\[
H_0:\mu_1-\mu_2=0,
\qquad
H_1:\mu_1-\mu_2\neq0.
\]

Bağımsızlık, veriye bakılarak sınanan sıradan bir dağılım varsayımı değildir;
örnekleme ve araştırma tasarımından anlaşılır. Aynı sınıftaki öğrenciler,
kardeş çiftleri veya bir danışmanın danışanları kümelenmiş olabilir. Böyle bir
yapıda sıradan bağımsız $t$ testi standart hatayı küçümseyebilir.

\subsection{Klasik havuzlanmış test ve Welch testi}
\index{t testi@$t$ testi!havuzlanmış}
\index{Welch t testi@Welch $t$ testi}
\index{varyans eşitliği}

Klasik bağımsız örneklem testi iki evren varyansının eşit olduğunu varsayar ve
birleşik varyans kullanır. Welch testi ise standart hatayı
\[
SE(\bar X_1-\bar X_2)
=\sqrt{\frac{s_1^2}{n_1}+\frac{s_2^2}{n_2}}
\]
ile hesaplar ve serbestlik derecesini varyanslarla grup büyüklüklerine göre
ayarlar:
\[
t=\frac{\bar X_1-\bar X_2}
{\sqrt{s_1^2/n_1+s_2^2/n_2}}.
\]

Welch testi eşit varyans varsayımı olmadan çalışır ve varyanslar eşit olduğunda
da genellikle çok az güç kaybeder. Bu nedenle özellikle grup büyüklükleri veya
standart sapmalar farklıysa güvenli bir varsayılan seçenektir
\cite{maxwell2018}.

\begin{mistakebox}
Levene testinde $p>0{,}05$ bulunması varyansların kesin eşit olduğunu
kanıtlamaz; küçük örneklemde testin gücü düşük olabilir. $p<0{,}05$ bulunması
da verinin analiz edilemeyeceği anlamına gelmez. Welch testi varyans eşitsizliği
için doğrudan bir çözüm sunar.
\end{mistakebox}

\subsection{Adım adım bağımsız grup örneği}

Program grubunda $n_1=40$, $\bar x_1=38{,}5$, $s_1=8{,}0$; karşılaştırma
grubunda $n_2=35$, $\bar x_2=44{,}0$, $s_2=10{,}0$ olsun. Fark ``program
eksi karşılaştırma'' olarak tanımlanmıştır.

\adimbaslik{1. Nokta tahmini}
\[
\bar x_1-\bar x_2=38{,}5-44{,}0=-5{,}5.
\]

\adimbaslik{2. Standart hata ve Welch istatistiği}
\[
SE=\sqrt{\frac{8^2}{40}+\frac{10^2}{35}}
\approx2{,}11,
\qquad
t=\frac{-5{,}5}{2{,}11}\approx-2{,}61.
\]

Welch serbestlik derecesi yaklaşık 65 ve çift yönlü $p$-değeri yaklaşık
0,011'dir.

\adimbaslik{3. Güven aralığı}

Yaklaşık kritik değer 2,00 alınırsa
\[
-5{,}5\pm2{,}00(2{,}11)
=[-9{,}72;\ -1{,}28]
\]
elde edilir. Veri program grubunda yaklaşık 1,3 ile 9,7 puan arasında daha
düşük evren ortalamasıyla uyumludur.

\adimbaslik{4. Standartlaştırılmış fark}

Birleşik standart sapma yaklaşık 8,99 olduğundan
\[
d=\frac{-5{,}5}{8{,}99}\approx-0{,}61.
\]
Negatif işaret program grubunun daha düşük puan aldığını, mutlak değer ise
farkın yaklaşık 0,61 standart sapma olduğunu gösterir.

\section{Eşleştirilmiş örneklem t testi}
\index{t testi@$t$ testi!eşleştirilmiş örneklem}
\index{fark puanı}

Eşleştirilmiş testte iki ölçüm aynı kişiye veya anlamlı biçimde eşleştirilmiş
birimlere aittir. Her birim için
\[
D_i=X_{i,\text{sonra}}-X_{i,\text{önce}}
\]
farkı hesaplanır. Hipotez ve test istatistiği
\[
H_0:\mu_D=0,
\qquad
t=\frac{\bar D}{s_D/\sqrt n},
\qquad sd=n-1
\]
biçimindedir. Dolayısıyla eşleştirilmiş $t$ testi, fark puanlarına uygulanan
tek örneklem $t$ testidir.

Eşleştirme, kişiler arası sabit farklılıkları fark hesabında ortadan kaldırır.
Ön-test ile son-test yüksek ilişkiliyse farkın standart hatası küçülebilir ve
aynı sayıda bağımsız gözleme göre daha hassas bir karşılaştırma elde edilir.

\subsection{Adım adım ön-test--son-test örneği}

Otuz öğrencinin son-test eksi ön-test farklarının ortalaması $\bar d=-4{,}5$
ve standart sapması $s_D=6{,}0$ olsun.

\[
SE(\bar D)=\frac{6}{\sqrt{30}}\approx1{,}10,
\qquad
t=\frac{-4{,}5}{1{,}10}\approx-4{,}11.
\]

Serbestlik derecesi 29 ve çift yönlü $p<0{,}001$'dir. Yüzde 95 güven aralığı
yaklaşık
\[
-4{,}5\pm2{,}045(1{,}10)
=[-6{,}74;\ -2{,}26]
\]
olur. Eşleştirilmiş standartlaştırılmış değişim
\[
d_z=\frac{\bar d}{s_D}=-0{,}75
\]
olarak hesaplanabilir.

\begin{sezgibox}
Eşleştirilmiş analiz, ``son-test grubunun ortalaması ön-test grubundan farklı
mı?'' sorusundan çok ``aynı öğrenciler tipik olarak ne kadar değişti?''
sorusunu yanıtlar. Analizin temel verisi iki ayrı sütun değil, her satırdaki
değişimdir.
\end{sezgibox}

\section{Eşleştirilmiş tasarımda dikkat edilmesi gerekenler}

Ön-test--son-test farkı bulunması değişimin yalnız programdan kaynaklandığını
göstermez. Zaman, olgunlaşma, sınanma etkisi, dış olaylar ve ortalamaya dönüş
değişime katkıda bulunabilir. Nedensel etki için uygun karşılaştırma grubu ve
tercihen rastgele atama gerekir.

Ayrıca:

\begin{itemize}
  \item Normallik iki ölçümün ayrı ayrı değil, fark puanlarının dağılımı için değerlendirilir.
  \item Bir öğrencinin ölçümlerinden biri eksikse o eşleşme klasik tam durum analizine giremez.
  \item Farkın yönü açıkça tanımlanmalıdır; ``sonra eksi önce'' ile ``önce eksi sonra'' işaretleri ters üretir.
  \item Üç veya daha fazla zaman noktası ayrı ayrı çok sayıda eşleştirilmiş testle incelenmemelidir.
  \item Gruplu ön-test--son-test araştırmasında yalnız her grubun kendi içindeki değişime bakmak, grupların değişim farkını doğrudan sınamaz.
\end{itemize}

Son madde özellikle önemlidir. Bir grupta $p<0{,}05$, diğerinde $p>0{,}05$
olması grupların değişimlerinin birbirinden anlamlı ölçüde farklı olduğunu
kanıtlamaz. Asıl hedef, değişim farkı veya uygun bir regresyon/ANCOVA modeliyle
grup etkisidir.

\begin{researchbox}
Ön-test puanları ölçüm hatası içeriyorsa yalnız uç puanlı öğrencileri seçmek
ortalamaya dönüş üretebilir. Son-testte daha ılımlı puan görülmesi program
etkisiyle karıştırılmamalıdır. Kontrol grubu ve ölçüm güvenirliği bu ayrımda
kritiktir.
\end{researchbox}

\section{Hangi t testi kullanılmalı?}

\begin{table}[htbp]
\centering
\caption{Araştırma durumuna göre test seçimi.}
\begin{tabular}{@{}p{0.28\textwidth}p{0.24\textwidth}p{0.27\textwidth}@{}}
\toprule
Durum & Uygun test & Sık yapılan hata \\
\midrule
Bir grubun ortalaması norm değerine karşı & Tek örneklem $t$ & Norm değerini örneklemden seçmek \\
Farklı öğrencilerden iki grup & Welch bağımsız $t$ & Grupları eşleştirilmiş saymak \\
Aynı öğrencilerin iki ölçümü & Eşleştirilmiş $t$ & Ölçümleri bağımsız analiz etmek \\
Üç bağımsız grup & ANOVA & Üç ayrı ikili test yürütmek \\
Üç veya daha fazla tekrarlı ölçüm & Tekrarlı ölçüm modeli & Çok sayıda eşleştirilmiş test yapmak \\
\bottomrule
\end{tabular}
\end{table}

Bir nicel sonuç ile iki kategorili bir grup değişkeni bulunması bağımsız $t$
testini otomatik olarak uygun yapmaz. Grup üyeliğinin gerçekten bağımsız olup
olmadığı ve sonuç değişkeninin anlamlı bir ortalamaya sahip olup olmadığı
değerlendirilmelidir.

\begin{etkinlik}
Şu durumların her biri için test seçin ve analiz birimini açıklayın: (a) bir
sınıfın kaygı ortalamasını norm değeriyle karşılaştırma, (b) kadın ve erkek
öğrencilerin uyku süresini karşılaştırma, (c) aynı danışanların ilk ve sekizinci
oturum puanlarını karşılaştırma, (d) üç terapi yaklaşımının ortalamalarını
karşılaştırma.
\end{etkinlik}

\section{Varsayımlar ve veri inceleme}
\index{varsayım kontrolü}
\index{bağımsızlık}
\index{normallik}

Üç $t$ testi için ortak ve tasarıma özgü kontroller vardır:

\begin{description}
  \item[Ölçme düzeyi] Sonuç değişkeninin ortalama ve fark işlemlerini anlamlı kılan nicel bir yapısı olmalıdır.
  \item[Bağımsızlık] Tek ve bağımsız testte gözlemler tasarıma uygun bağımsızdır; eşleştirilmiş testte çiftler birbirinden bağımsızdır.
  \item[Dağılım biçimi] Küçük örneklemde güçlü çarpıklık ve ağır kuyruklar daha önemlidir.
  \item[Aykırı değer] Ortalama ve standart sapmayı etkileyerek test sonucunu değiştirebilir.
  \item[Varyans yapısı] Klasik bağımsız test eşit varyans varsayar; Welch testi bu koşulu gerektirmez.
  \item[Örnekleme] Uygun olmayan seçim süreci, doğru hesaplanan testi hedef evrene genellenemez hâle getirebilir.
\end{description}

Histogram, kutu grafiği ve Q--Q grafiği sayısal özetlerle birlikte incelenir.
Küçük örneklemde tek bir aykırı değer sonucu güçlü biçimde etkileyebilir. Büyük
örneklemde $t$ testi orta düzey dağılım sapmalarına daha dayanıklıdır; ancak
bağımsızlık ihlaline veya sistematik yanlılığa dayanıklı değildir.

\begin{assumptionbox}
Normallik testi anlamlı değil diye dağılımın normal olduğu kanıtlanmış sayılmaz;
anlamlı çıktı diye de otomatik olarak parametrik olmayan teste geçilmez.
Örneklem hacmini, sapmanın biçimini, aykırı değerleri ve $t$ testinin
sağlamlığını birlikte değerlendirin.
\end{assumptionbox}

\section{Aykırı değer ve duyarlılık analizi}
\index{aykırı değer}
\index{duyarlılık analizi}

Aykırı bir puan önce veri kaydı, puanlama ve katılımcı bağlamı bakımından
incelenmelidir. Gerçek bir gözlem gerekçesiz biçimde silinmemelidir. Sonuç uç
gözleme duyarlıysa:

\begin{itemize}
  \item gözlem dahil ve hariç sonuçlar karşılaştırılabilir,
  \item ham fark yanında medyan ve IQR verilebilir,
  \item sağlam veya yeniden örneklemeye dayalı aralıklar düşünülebilir,
  \item uygun parametrik olmayan yöntemle sonuçların yönü incelenebilir.
\end{itemize}

Duyarlılık analizi, istenen $p$-değerini veren seçeneği seçme yolu değildir.
Ana analiz önceden belirlenmeli, alternatif sonuçlar şeffaf biçimde
raporlanmalıdır.

\section{Parametrik olmayan seçenekler}
\index{parametrik olmayan yöntem}

Tek örneklem veya eşleştirilmiş durumda işaret testi ya da Wilcoxon işaretli
sıra testi; bağımsız iki grupta Mann--Whitney testi düşünülebilir. Bu testler
``ortalamaların parametrik olmayan $t$ testi'' olarak görülmemelidir.
Yorumları dağılım, sıralama ve bazı ek koşullar altında konum farkına dayanır.

Yöntem seçerken şu sorular sorulur:

\begin{itemize}
  \item Araştırmanın hedefi ortalama farkı mı, medyan veya sıralama davranışı mı?
  \item Dağılım sorunu birkaç veri hatasından mı, gerçek çarpıklıktan mı kaynaklanıyor?
  \item Örneklem hacmi ve $t$ testinin sağlamlığı yeterli mi?
  \item Sağlam ortalama, permütasyon veya bootstrap yöntemi daha doğrudan bir hedef sunuyor mu?
\end{itemize}

Parametrik olmayan yöntemler Bölüm 13'te ayrıntılı olarak ele alınacaktır.

\section{Etki büyüklüğü ve ham fark}
\index{etki büyüklüğü}
\index{Cohen d@Cohen'in $d$ değeri}
\index{ham ortalama farkı}

Ham ortalama farkı ölçme aracının özgün biriminde doğrudan yorumlanır.
Standartlaştırılmış fark ise farklı ölçek veya çalışmalar arasında
karşılaştırmayı kolaylaştırır. Bağımsız gruplarda yaygın ölçü
\[
d=\frac{\bar x_1-\bar x_2}{s_{\text{birleşik}}}
\]
biçimindedir. Küçük örneklemde Cohen'in $d$ değerinin yanlılığını azaltan
Hedges'in $g$ değeri kullanılabilir.

Eşleştirilmiş tasarımda standartlaştırma için fark puanlarının standart
sapmasını kullanan $d_z=\bar d/s_D$ yaygındır. Ancak bağımsız ve eşleştirilmiş
tasarımlardaki $d$ ölçüleri farklı paydalara dayanabilir; raporda kullanılan
formül belirtilmelidir.

\begin{mistakebox}
``$d=0{,}50$ orta etkidir'' etiketi bağlamın yerine geçmez. Beş puanlık fark
bir ölçekte önemsiz, başka bir ölçekte hizmet kararını değiştirecek kadar
önemli olabilir. Ham fark, güven aralığı ve uygulamada anlamlı en küçük değişim
birlikte verilmelidir.
\end{mistakebox}

\section{Örneklem hacmi ve güç}
\index{örneklem hacmi}
\index{test gücü}

Bir $t$ testinin gücü gerçek etki büyüklüğü, örneklem hacmi, değişkenlik,
anlamlılık düzeyi ve tasarımdan etkilenir. Bağımsız iki grupta toplam örneklem
aynı olsa bile grupların dengeli olması çoğunlukla daha yüksek hassasiyet
sağlar. Eşleştirilmiş tasarımda ölçümler arası ilişki arttıkça fark puanının
değişkenliği azalabilir.

Örneklem hacmi ``en az 30 kişi'' gibi evrensel bir kuralla belirlenmemelidir.
Planlamada:

\begin{itemize}
  \item uygulamada anlamlı en küçük fark,
  \item beklenen standart sapma,
  \item hedef güç ve $\alpha$ düzeyi,
  \item tek veya çift yönlü test,
  \item grup oranı, kümelenme ve beklenen veri kaybı
\end{itemize}

birlikte değerlendirilir. Etki büyüklüğünü önceki küçük çalışmanın iyimser
sonucundan doğrudan almak yetersiz örneklem planına yol açabilir.

\begin{researchbox}
Etik açıdan gereğinden küçük örneklem, katılımcı yüküne rağmen araştırma
sorusunu yanıtlayamayabilir. Gereğinden büyük örneklem de kaynak ve katılımcı
zamanı tüketebilir. Güç analizi yalnız teknik değil, etik bir planlama
aracıdır.
\end{researchbox}

\section{Python ile üç t testi}

Aşağıdaki kod özet istatistiklerden tek örneklem ve bağımsız Welch testini,
ham fark verisinden eşleştirilmiş testi yeniden üretir:

\begin{lstlisting}
import numpy as np
from scipy import stats

# 1) Tek orneklem: xbar=54.8, s=8, n=25, mu0=50
n, xbar, s, mu0 = 25, 54.8, 8.0, 50
se = s / np.sqrt(n)
t_tek = (xbar - mu0) / se
p_tek = 2 * stats.t.sf(abs(t_tek), df=n - 1)
kritik = stats.t.ppf(0.975, df=n - 1)
ga_tek = (xbar - kritik * se, xbar + kritik * se)

# 2) Bagimsiz gruplar: ozetlerden Welch testi
welch = stats.ttest_ind_from_stats(
    mean1=38.5, std1=8.0, nobs1=40,
    mean2=44.0, std2=10.0, nobs2=35,
    equal_var=False
)

# 3) Eslestirilmis test: her ogrencinin sonra-once farki
farklar = np.array([-8, -3, -5, 2, -6, -4, -9, -1,
                    -7, -2, -5, -6, -3, -8, 0, -4])
esli = stats.ttest_1samp(farklar, popmean=0)

print("Tek orneklem:", t_tek, p_tek, ga_tek)
print("Welch testi:", welch)
print("Eslestirilmis test:", esli)
print("Ortalama degisim:", farklar.mean())
\end{lstlisting}

Gerçek veri analizinde koddan önce değişken tanımları, eksik gözlemler,
olanaksız değerler ve grafikler incelenmelidir. Sonuçlar yazılım nesnesinden
alınırken farkın yönü ve varsayılan alternatif hipotez doğrulanmalıdır.

\section{SPSS ile uygulama ve çıktı okuma}

\subsection*{Tek örneklem t testi}

\begin{enumerate}
  \item \textbf{Analyze $>$ Compare Means $>$ One-Sample T Test} yolunu izleyin.
  \item Nicel değişkeni \textbf{Test Variable(s)} alanına taşıyın.
  \item \textbf{Test Value} alanına kuramsal veya norm değerini girin.
  \item Çıktıda ortalama, standart sapma, $t$, $sd$, $p$, ortalama farkı ve fark güven aralığını okuyun.
\end{enumerate}

\subsection*{Bağımsız örneklem t testi}

\begin{enumerate}
  \item \textbf{Independent-Samples T Test} menüsünü açın.
  \item Sonuç değişkenini test alanına, grup değişkenini gruplama alanına taşıyın.
  \item \textbf{Define Groups} ile iki grup kodunu doğru tanımlayın.
  \item Grup özetlerini ve farkın hangi grup eksi hangi grup olduğunu kontrol edin.
  \item Varyans yapısına uygun satırdan, tercihen Welch sonucundan, test ve güven aralığını okuyun.
\end{enumerate}

\subsection*{Eşleştirilmiş örneklem t testi}

\begin{enumerate}
  \item \textbf{Paired-Samples T Test} menüsünde ön-test ve son-test değişkenlerini bir çift olarak tanımlayın.
  \item Değişkenlerin sıraya göre hangi farkı ürettiğini doğrulayın.
  \item Eşleşmiş gözlem sayısını ve eksik çiftleri kontrol edin.
  \item Fark ortalaması, fark standart sapması, $t$, $sd$, $p$ ve güven aralığını birlikte okuyun.
\end{enumerate}

\begin{outputguidebox}
SPSS çıktısını şu sırayla okuyun: (1) geçerli grup veya çift sayıları, (2)
ortalamalar ve standart sapmalar, (3) farkın yönü, (4) standart hata, (5)
$t$, serbestlik derecesi ve kesin $p$-değeri, (6) farkın güven aralığı, (7)
ayrıca hesaplanan etki büyüklüğü. Levene tablosunu ana araştırma sonucu gibi
raporlamayın.
\end{outputguidebox}

\section{Üç örnek raporlama kalıbı}

\subsection*{Tek örneklem}

``Öğrencilerin akademik kaygı ortalaması ($\bar x=54{,}8$, $SS=8{,}0$,
$n=25$), norm değeri 50'den yüksektir, $t(24)=3{,}00$, $p=0{,}006$,
yüzde 95 GA $[51{,}50;58{,}10]$, $d=0{,}60$.''

\subsection*{Bağımsız örneklem}

``Program grubunun yalnızlık ortalaması ($\bar x=38{,}5$, $SS=8{,}0$,
$n=40$), karşılaştırma grubundan ($\bar x=44{,}0$, $SS=10{,}0$, $n=35$)
5,5 puan daha düşüktür. Welch testi fark için kanıt göstermiştir,
$t(65)=-2{,}61$, $p=0{,}011$, yüzde 95 GA $[-9{,}72;-1{,}28]$,
$d=-0{,}61$.''

\subsection*{Eşleştirilmiş örneklem}

``Otuz öğrencinin son-test puanı ön-testten ortalama 4,5 puan daha düşüktür
($SS_{fark}=6{,}0$). Değişim istatistiksel olarak belirgindir,
$t(29)=-4{,}11$, $p<0{,}001$, yüzde 95 GA $[-6{,}74;-2{,}26]$,
$d_z=-0{,}75$.''

Bu kalıplar tasarımın nedensel yoruma izin verdiği anlamına gelmez. Atama,
örnekleme ve kayıp katılımcı süreçleri yöntem ve sınırlılıklar bölümünde
ayrıca açıklanmalıdır.

\section{Bir t testi araştırmasını değerlendirirken}

\begin{enumerate}
  \item Karşılaştırma sabit değerle mi, bağımsız gruplar arasında mı, eşleşmiş ölçümler arasında mı yapılmıştır?
  \item Hedef farkın yönü ve sıfır hipotezi açık mıdır?
  \item Bağımsızlık yapısı araştırma tasarımıyla uyumlu mudur?
  \item Dağılım, aykırı değer, varyans ve eksik veri sorunları incelenmiş midir?
  \item Bağımsız gruplarda Welch yaklaşımı veya eşit varyans gerekçesi belirtilmiş midir?
  \item Ham fark, güven aralığı ve uygun etki büyüklüğü $p$-değeriyle birlikte verilmiş midir?
  \item Sonuç araştırma tasarımının izin verdiğinden daha güçlü bir nedensel ya da klinik dille sunulmuş mudur?
\end{enumerate}

\section{R ile çözümlü uygulama}
\label{sec:r-b11}

\textbf{Soru.} Bölümdeki bağımsız grupların özetlerinden Welch testini;
yirmi öğrencinin fark puanlarının özetinden eşleştirilmiş testi hesaplayın.
Ham veriler verilmediğinden aşağıdaki çözümler özet istatistik formüllerini kullanır.

\begin{lstlisting}[style=prismR]
hacim <- c(30, 28)
ortalama <- c(78, 72)
standart_sapma <- c(8, 10)
katki <- standart_sapma^2 / hacim
standart_hata <- sqrt(sum(katki))
fark <- ortalama[1] - ortalama[2]
serbestlik <- sum(katki)^2 / sum(katki^2 / (hacim - 1))
t_degeri <- fark / standart_hata
c(t = t_degeri, sd = serbestlik,
  p = 2 * pt(abs(t_degeri), df = serbestlik,
             lower.tail = FALSE))
fark + c(-1, 1) * qt(0.975, df = serbestlik) * standart_hata
cift_sayisi <- 20
ortalama_fark <- 4.2
fark_sapmasi <- 5
es_se <- fark_sapmasi / sqrt(cift_sayisi)
es_t <- ortalama_fark / es_se
c(t = es_t, sd = cift_sayisi - 1,
  p = 2 * pt(abs(es_t), df = cift_sayisi - 1,
             lower.tail = FALSE),
  dz = ortalama_fark / fark_sapmasi)
ortalama_fark + c(-1, 1) *
  qt(0.975, df = cift_sayisi - 1) * es_se
\end{lstlisting}

\begin{outputguidebox}
\textbf{Çözüm ve kontrol değerleri.} Welch testinde fark 6 puan,
$t=2{,}5121$, serbestlik derecesi 51,7113, $p=0{,}01516$ ve yüzde 95
fark aralığı $[1{,}2066;10{,}7934]$'tür. Eşleştirilmiş testte
$t(19)=3{,}7566$, $p=0{,}001336$, aralık $[1{,}8599;6{,}5401]$ ve
$d_z=0{,}84$ bulunur.
\end{outputguidebox}

\textbf{Yorum.} Ham veriyle \texttt{t.test} kullanılırken bağımsız gruplar
için \texttt{var.equal = FALSE}, eşleşmiş son ve ön ölçüm vektörleri için
\texttt{paired = TRUE} seçilir. Eşleştirilmiş testte aynı
konumdaki iki değer aynı kişiye ait olmalıdır. Özetleri ham gözlem gibi bu
komutlara vermeyin; 51,7113 serbestlik derecesini de tam sayıya yuvarlamayın.

\textbf{Pekiştirme ve yanıt.} Fark tanımı ``ön--son'' yapılırsa eşleştirilmiş
$t$ ve $d_z$ işaret değiştirir; çift yönlü \pval-değeri değişmez ve güven
aralığı $[-6{,}5401;-1{,}8599]$ olur.
\section{Bölüm sonu çözümlü problemler}

\begin{enumerate}
  \item \textbf{Tek örneklem testi.} $n=16$, $\bar x=54$, $s=8$ ve $\mu_0=50$ için $t$ değerini ve serbestlik derecesini bulun.

  \textbf{Çözüm:} $SE=8/\sqrt{16}=2$ ve $t=(54-50)/2=2$'dir. Serbestlik derecesi $16-1=15$ olur. Karar için alternatif hipotez ve $t(15)$ dağılımından $p$-değeri gerekir.

  \item \textbf{Welch testi.} İki grubun ortalamaları 40 ve 46; standart sapmaları 6 ve 8; örneklem hacimleri 25 ve 16'dır. Farkı ve Welch standart hatasını bulun.

  \textbf{Çözüm:} Birinci eksi ikinci farkı $40-46=-6$'dır. $SE=\sqrt{6^2/25+8^2/16}=\sqrt{1{,}44+4}=\sqrt{5{,}44}\approx2{,}33$ ve $t\approx-6/2{,}33=-2{,}57$ olur.

  \item \textbf{Eşleştirilmiş test.} On öğrencinin son-test eksi ön-test fark ortalaması $-3$, fark standart sapması 4'tür. Test istatistiğini ve $d_z$ değerini bulun.

  \textbf{Çözüm:} $SE=4/\sqrt{10}\approx1{,}265$ ve $t=-3/1{,}265\approx-2{,}37$'dir; $sd=9$ olur. $d_z=\bar d/s_D=-3/4=-0{,}75$'tir. Negatif yön, son-test puanlarının daha düşük olduğunu gösterir.
\end{enumerate}

\bolumkaynaklari{\cite{casella2002,maxwell2018,cohen1988,wasserman2004}}

\section*{Hatırla--Uygula--Yansıt}

\begin{enumerate}
\renewcommand{\labelenumi}{\textbf{\Alph{enumi}.}}
  \item \textbf{Hatırla:} Tek örneklem, bağımsız örneklem ve eşleştirilmiş örneklem $t$ testlerinin hedef parametrelerini yazın.
  \item \textbf{Seç:} Bir sınıfın ortalamasını normla, iki farklı sınıfı ve aynı öğrencilerin iki ölçümünü karşılaştırmak için uygun testleri belirleyin.
  \item \textbf{Hesapla:} $\bar x=62$, $s=10$, $n=25$ ve $\mu_0=58$ için tek örneklem $t$ istatistiğini bulun.
  \item \textbf{Yorumla:} Tek örneklem farkı için yüzde 95 güven aralığı $[1{,}2;6{,}8]$ olduğunda sonucu açıklayın.
  \item \textbf{Welch:} $n_1=30$, $s_1=6$, $n_2=45$, $s_2=10$ için ortalama farkının standart hatasını hesaplayın.
  \item \textbf{Eşleştir:} $D=$ sonra eksi önce olarak tanımlandığında $\bar D=-3$ ne anlama gelir? Fark yönü ters seçilirse ne değişir?
  \item \textbf{Varsayım:} Bağımsızlık varsayımının neden Shapiro--Wilk veya Levene testiyle sınanamayacağını açıklayın.
  \item \textbf{Eleştir:} ``Levene testi anlamlı olmadığı için varyanslar kesinlikle eşittir'' ifadesini düzeltin.
  \item \textbf{Etki büyüklüğü:} Ortalama fark 7 ve birleşik standart sapma 14 ise Cohen'in $d$ değerini hesaplayıp bağlamla yorumlayın.
  \item \textbf{Ön--son tasarım:} Bir grupta değişim anlamlı, diğerinde anlamsız bulunduğunda neden grupların değişimlerinin farklı olduğu sonucuna doğrudan varılamaz?
  \item \textbf{Duyarlılık:} Tek güçlü aykırı değerin grup ortalaması, standart sapma, $t$ ve $d$ üzerindeki olası etkilerini tartışın.
  \item \textbf{Uygula:} Python kodunda grup standart sapmalarını ve örneklem sayılarını değiştirerek Welch sonucunu inceleyin.
  \item \textbf{Raporla:} İki bağımsız grup için $n$, ortalama, standart sapma, fark, $t$, $sd$, $p$, güven aralığı ve $d$ içeren kısa bir bulgu paragrafı yazın.
  \item \textbf{Yansıt:} PDR programında istatistiksel olarak anlamlı ortalama azalmanın neden her öğrencinin iyileştiği anlamına gelmediğini açıklayın.
\end{enumerate}